\chapter{Introducción}

\section{Planteamiento del problema}

La detección de anomalías en astronomía de dominio temporal es clave para identificar fenómenos poco frecuentes o desconocidos. La operación del Observatorio Vera C. Rubin incrementará de forma sustantiva el volumen y la complejidad de los datos fotométricos multibanda, caracterizados por muestreo irregular, relaciones entre bandas y gran variabilidad intrínseca. Este contexto plantea dos desafíos complementarios: (i) seleccionar de manera eficiente objetos potencialmente interesantes dentro de un flujo masivo de alertas y (ii) comprender, con criterios reproducibles, por qué dichos objetos se apartan del comportamiento esperado.

No basta con marcar un ejemplo como atípico; es necesario explicarlo. Entre las aproximaciones de interpretabilidad, las explicaciones contrafactuales ofrecen un marco explícito: describen modificaciones mínimas y plausibles de los datos que harían que el ejemplo dejara de ser considerado anómalo. Estas explicaciones facilitan la priorización de seguimientos y la contrastación de hipótesis, al vincular la detección con una narrativa concreta de cambio en las curvas de luz.

\section{Justificación}

El dataset ELAsTiCC-1, desarrollado en el contexto del Observatorio Vera C. Rubin, contiene cerca de 1.4 millones de curvas de luz simuladas de 17 tipos de supernovas. La escala y complejidad de este dataset hacen inviable la inspección manual, y los métodos de clasificación supervisada requieren etiquetas que en la práctica no siempre estarán disponibles en tiempo real. Un sistema de detección de anomalías no supervisado, capaz además de generar explicaciones interpretables, es fundamental para el flujo de trabajo de la astronomía de grandes volúmenes de datos.

Desde el punto de vista metodológico, la representación de curvas de luz multibanda como imágenes 2grid —donde las seis bandas fotométricas se organizan como subgráficos apilados— permite aprovechar arquitecturas convolucionales estándar y capturar de forma natural la correlación temporal entre bandas, eliminando la necesidad de manejar muestreo irregular.

\section{Objetivos}

\subsection{Objetivo general}

Desarrollar un pipeline de detección de anomalías e interpretación contrafactual para curvas de luz de supernovas del dataset ELAsTiCC-1, basado en una representación discreta del espacio latente aprendida con un autoencoder con cuantización vectorial (VQ-VAE) y un modelo prior tipo MaskGIT.

\subsection{Objetivos específicos}

\begin{enumerate}
    \item Diseñar e implementar una representación 2grid de curvas de luz multibanda como imágenes RGB de 256×256 píxeles, con normalización por objeto que garantice comparabilidad entre supernovas de distinta luminosidad.

    \item Entrenar un VQ-VAE con codebook EMA y mecanismo de codebook reset sobre el dataset ELAsTiCC-1, evaluando la calidad de reconstrucción y la utilización del espacio latente discreto.

    \item Estudiar el efecto de agregar clasificación supervisada sobre el espacio cuantizado $\mathbf{z}_q$, evaluando si la organización semántica inducida en el codebook mejora la interpretabilidad de los contrafactuales generados.

    \item Entrenar un prior tipo MaskGIT sobre las secuencias de tokens del VQ-VAE y utilizarlo para asignar scores de anomalía y generar contrafactuales mediante completado condicional de tokens.

    \item Evaluar el pipeline de detección de anomalías y la calidad de los contrafactuales generados según criterios de plausibilidad, parsimonia y consistencia temporal entre bandas.
\end{enumerate}
